\chapter{Ionized Calcium}

\section*{What Is This Test?}
Ionized calcium measures the free, biologically active fraction of calcium in blood. Unlike total calcium, it is less dependent on albumin concentration but is sensitive to pH and specimen handling.

\section*{Alternative Names}
Free Calcium; Serum Ionized Calcium; Ca$^{2+}$.

\section*{Why This Test Is Ordered}
It is particularly useful in critically ill patients, patients with abnormal albumin or acid-base status, suspected severe hypercalcemia or hypocalcemia, disorders of parathyroid function, massive transfusion, pancreatitis, and unexplained neuromuscular or cardiac symptoms.

\section*{How This Test Is Performed}
Blood is collected anaerobically into a suitable heparinized device and analyzed promptly on a blood-gas or chemistry analyzer. The laboratory records or corrects for pH because alkalosis lowers the ionized fraction without necessarily changing total calcium.

\section*{Preparation, Risks, and Limitations}
No fasting is usually needed. The main risks are those of venipuncture. Air exposure, delays, excess liquid heparin, temperature changes, and abnormal pH can produce inaccurate results. Ionized calcium should not be interpreted as interchangeable with total calcium.

\section*{Understanding Results}
Low ionized calcium may cause tingling, muscle spasm, seizures, or QT prolongation; high values may cause weakness, confusion, constipation, or rhythm disturbances. Interpretation requires albumin, phosphate, magnesium, PTH, kidney function, pH, and the clinical condition.

\section*{References}
International Federation of Clinical Chemistry guidance on ionized calcium measurement.

\clearpage
\section*{Understanding Results and Follow-up}
The laboratory report should identify the specimen, method, units, reference interval or decision limit, and whether the result is positive, negative, abnormal, or inconclusive. Reference intervals vary with age, sex, pregnancy, specimen type, assay, and laboratory; a result outside a range is not automatically a diagnosis, and a result within a range does not always exclude disease. Discuss unexpected results with the ordering clinician, who can decide whether repeat or confirmatory testing, treatment, or referral is appropriate. Do not change medicines or delay urgent care based on an isolated result.


\section*{Regulatory Status and Authorization Date}
Regulatory status applies to the specific assay, instrument, manufacturer, and intended use rather than to the test name alone. No single approval date can be assigned to this test category. For a commercial assay, verify the current product in the relevant regulator's database: the U.S. Food and Drug Administration (FDA) clearance, approval, or authorization; India's Central Drugs Standard Control Organisation (CDSCO) licence or permission; the European Union In Vitro Diagnostic Regulation (EU IVDR) CE certificate issued through the applicable conformity-assessment route; or the United Kingdom Medicines and Healthcare products Regulatory Agency (MHRA) registration and UKCA/CE status. Laboratory-developed tests may not have an individual FDA, CDSCO, EU IVDR, or MHRA approval date.
\clearpage
